\documentclass{article}
\usepackage{amsmath, amssymb}
\usepackage[russian]{babel}

\DeclareMathOperator{\Tr}{Tr}
\DeclareMathOperator{\dom}{dom}
\DeclareMathOperator{\KL}{KL}

\begin{document}

\section*{Выпуклые функции}

\section*{Задача 1}

\subsection*{1.1}
Функция $ f(x, y) = x^4 + x^2 y^2 + y^4 $ является выпуклой.

$$
H(x,y) = \begin{bmatrix}
12x^2 + 2y^2 & 4xy \\
4xy & 2x^2 + 12y^2
\end{bmatrix}.
$$
Определитель гессиана: 
$$
\det(H) = 24x^4 + 132x^2y^2 + 24y^4> 0
$$

так же $12x^2 + 2y^2>0$ и $12y^2 + 2x^2 > 0$ $\rightarrow$ гессиан положительно определен. Тогда функция выпукла.

\subsection*{1.2}
Функция $f(x, y) = x^2 y^2$ не является выпуклой.
$$
H(x,y) = \begin{bmatrix}
2y^2 & 4xy \\
4xy & 2x^2
\end{bmatrix}.
$$
В точке $(1,1)$ гессиан равен:

$det(H)(1,1) = -12$ следовательно, функция не выпукла.

\subsection*{1.3}
Функция  $f(x, y) = \|x - y\| + \|x + y\| $ выпукла.\\
Функции $g(x,y) = \|x - y\| $ и $h(x,y) = \|x + y\|$  являются выпуклыми, так как норма — выпуклая функция, а аффинные преобразования сохраняют выпуклость. Сумма выпуклых функций выпукла.

\subsection*{1.4}
Функция $f(x) = \tan ( \|x\|_\infty^3 ) $ с областью определения  $\dom f = [0,1]^d $ является выпуклой. \\
Функция $\|x\|_\infty$  выпукла на  $\mathbb{R}^d $. Функция  $t^3$  выпукла и возрастает на $[0,1]$, поэтому композиция  $\|x\|_\infty^3$  выпукла. Функция  $\tan(t)$  выпукла и возрастает на $[0,1]$ Таким образом, композиция  $\tan(\|x\|_\infty^3)$  выпукла.

\subsection*{1.5}
Функция \( f(x) = \sum_{i=1}^d x_i^4 \) выпукла, но не является $\mu$-сильно выпуклой ни для какого $\mu > 0$. \\

Каждая функция $x_i^4$ выпукла на $\mathbb{R}$, поэтому их сумма выпукла. Однако гессиан $f$ равен диагональной матрице с элементами $(12x_i^2)$. В точке $x = 0$ гессиан нулевой, поэтому не  $\mu$ сильно выпукла.


\subsection*{2} Докажите, что KL-дивергенция \( f : \mathbb{R}^d \to \mathbb{R} \):
\[
f(p) = \KL(p \mid q) = \sum_{i=1}^d p_i \log \frac{p_i}{q_i}, \quad \dom f = \left\{ p \in \mathbb{R}_+^d \middle| \sum_{i=1}^d p_i = 1 \right\},
\]
где \( q \in \left\{ q \in \mathbb{R}_+^d \middle| \sum_{i=1}^d q_i = 1 \right\} \) является неотрицательной.

Указание: воспользуйтесь неравенством Йенсена.

\subsection*{Решение}

Заметим, что KL-дивергенцию можно переписать в виде:
$$
\KL(p \mid q) = \sum_{i=1}^d p_i \log \frac{p_i}{q_i} = -\sum_{i=1}^d p_i \log \frac{q_i}{p_i}.
$$

Рассмотрим функцию  $-\log t $. Эта функция выпукла на  $\mathbb{R}_{++} $

Применим неравенство Йенсена к выпуклой функции


$$
\sum_{i=1}^d p_i * f(x_i) > f ( \sum_{i=1}^d p_i * x_i )\quad \text{где } \quad \sum_{i=1}^d p_i = 1
$$
Подставим $-\log t$ вместо $f$
$$
\sum_{i=1}^d p_i * -\log \frac{q_i}{p_i} > -\log  ( \sum_{i=1}^d p_i * \frac{q_i}{p_i} ) = 0
$$

Таким образом, $\KL(p \mid q) \geq 0$  для всех  $p$ и $q$. Равенство достигается тогда и только тогда, когда $ p = q $


\subsection*{3} Пусть \( g : \mathbb{R}_+ \to \mathbb{R}_+ \) — непрерывная выпуклая функция. Определим \( f : \mathbb{R}_+ \to \mathbb{R}_+ \):
\[
f(x) = \frac{1}{x} \int_0^x g(t) \, dt.
\]
Покажите, что \( f \) тоже выпукла.

\subsection*{Решение}

пусть $t = xr \rightarrow dt = xdr$

$$
\frac{1}{x} \int_0^x g(t) \, dt = \frac{1}{x} \int_0^1 g(xr) x dr = \int_0^1 g(xr)dr
$$

Проверим выпуклость по определению


$$
\int_0^1 g((\alpha x_1)r)dr + \int_0^1 g(((1 -\alpha) x_2)r)dr = 
$$

$$
= \int_0^1 (g((\alpha x_1)r) +  g(((1 -\alpha) x_2)r))dr > \int_0^1 (g((\alpha x_1 + (1 - \alpha)r) =
$$

$$
f(\alpha x_1 + (1 - \alpha)x_2)
$$

что и требовалось доказать



\subsection*{4} Проверьте, является ли выпуклой функция \( f : \mathbb{S}_{++}^d \to \mathbb{R} \):
\[
f(X) = \Tr(X^{-1}).
\]

\textbf{Решение}

$$
Df[H] = D(Tr(X^{-1})) = D<X^{-1}, I>[H] = <-X^{-1}HX^{-1}, I> = -<X^{-1}HX^{-1}, I>
$$

$$
D^2f[H] =2[<X^{-1}HX^{-1} HX^{-1}, I>]
$$

Заметим что

$$
Tr(X^{-1}VX^{-1}VX^{-1}) = Tr((X^{-1/2}VX^{-1/2})^2 X^{-1})
$$

$X^{-1}$ положительно опрелеленная функция. Получается $Tr((X^{-1/2}VX^{-1/2})^2 X^{-1}) > 0$

Значит f - выпуклая


\subsection*{5}

Используя неравенство Йенсена, докажите неравенство Юнга:

\[
|xy| \leq \frac{|x|^p}{p} + \frac{|y|^q}{q},
\]
где $\frac{1}{p} + \frac{1}{q} = 1$

\subsection*{Решение}

Рассмотрим $e^t$, которая выпукла на $R$

Применим неравенство Йенсена к выпуклой функции с весами $\frac{1}{p}$ и $\frac{1}{q}$:

$$
\frac{1}{p} e^{t_1} + \frac{1}{q} e^{t_2} \geq e^{\frac{t_1}{p} + \frac{t_2}{q}}.
$$

Теперь положим:
$$
t_1 = p \ln |x|, \quad t_2 = q \ln |y|.
$$

Тогда:
$$
\frac{1}{p} e^{p \ln |x|} + \frac{1}{q} e^{q \ln |y|} \geq e^{\frac{p \ln |x|}{p} + \frac{q \ln |y|}{q}} = e^{\ln |x| + \ln |y|} = e^{\ln |xy|} = |xy|.
$$

Упрощаем левую часть:
$$
\frac{1}{p} e^{p \ln |x|} = \frac{1}{p} |x|^p, \frac{1}{q} e^{q \ln |y|} = \frac{1}{q} |y|^q.
$$

Таким образом, получаем:
$$
\frac{|x|^p}{p} + \frac{|y|^q}{q} \geq |xy|,
$$
что и требовалось доказать.


\textbf{Задача 6.} (1 балл) Проверьте, являются ли выпуклыми функции \( f : \mathbb{S}^d \to \mathbb{R} \):  
Указание: воспользуйтесь соотношением Рэлея.

\textbf{6.1.} (0.5 балла)
\[
f(X) = \lambda_{\max}(X).
\]

\textbf{6.2.} (0.5 балла)
\[
f(X) = \max \left\{ |\lambda_{\min}(X)|, |\lambda_{\max}(X)| \right\}.
\]

\end{document}